\rtmtight
\begin{tblr}{width=\textwidth,colspec={|Q[c,m,2.1cm]|X[l,m]|},hlines,vlines,hspan=minimal,rowsep=1.5pt,colsep=3pt,row{1}={bg=headgray,font=\bfseries}}
\SetCell{c,m}\hfil\logo\hfil & \textbf{POLITEKNIK NEGERI MALANG}\par\textbf{JURUSAN TEKNOLOGI INFORMASI}\par\textbf{PROGRAM STUDI: D4 TEKNIK INFORMATIKA} \\
\SetCell[c=2]{l} \textbf{RENCANA TUGAS MAHASISWA (RTM) DAN RUBRIK PENILAIAN} & \\
Nama Mata Kuliah & Fisika \\
Kode Mata Kuliah & RTI251009 \quad \textbf{sks / Jam perminggu:} 2 SKS/2 jam \quad \textbf{Semester:} 1 \\
Dosen Pengampu & \begin{enumerate}\item Vipkas Al Hadid Firdaus, S.T., M.T.\item Sofyan Noor Arief, S.ST., M.Kom.\item Agung Nugroho Pramudhita, S.T., M.T.\end{enumerate} \\
\end{tblr}
\assessmentblock{Tugas Terstruktur Penyelesaian Masalah Fisika}{Tugas terstruktur}{SCPMK0902-00901--SCPMK0902-00904}{Mahasiswa menyelesaikan tugas terstruktur pemecahan masalah fisika pada setiap tahapan materi (pengukuran, mekanika, energi-momentum, dan kelistrikan) dengan konteks sistem teknologi informasi.}{Mengerjakan soal dan studi kasus secara terstruktur: menuliskan diketahui-ditanya, memilih hukum/persamaan yang sesuai, menghitung, memeriksa satuan, dan menginterpretasikan hasil pada konteks sistem.}{Lembar jawaban atau laporan tugas tertulis/digital yang memuat langkah penyelesaian lengkap.}{Penyelesaian tugas konsep dasar, metode ilmiah, dan pengukuran & 5 & Jawaban tugas pekan 1--3 menerapkan konsep besaran, satuan, dan ketidakpastian dengan konversi benar serta langkah pengukuran dijelaskan runtut. & Konsep dasar sebagian besar benar, tetapi ada kesalahan konversi satuan atau perhitungan ketidakpastian. & Jawaban tidak menggunakan konsep besaran/satuan yang benar atau langkah pengukuran tidak dapat diikuti. \\ \\
Penyelesaian tugas kinematika dan dinamika partikel & 6 & Persamaan gerak dan hukum Newton dipilih tepat, diagram gaya digambar benar, dan hasil hitung percepatan/gaya benar dengan satuan konsisten. & Pemilihan persamaan umumnya tepat, tetapi ada kesalahan pada diagram gaya atau perhitungan. & Persamaan gerak atau hukum Newton diterapkan keliru sehingga hasil tidak benar. \\ \\
Penyelesaian tugas usaha, energi, dan momentum & 3 & Hukum kekekalan energi dan momentum diterapkan tepat pada kasus konversi energi dan tumbukan dengan langkah dan hasil benar. & Konsep kekekalan digunakan, tetapi ada kesalahan perhitungan atau identifikasi jenis tumbukan. & Hukum kekekalan tidak diterapkan atau hasil tidak relevan dengan kasus. \\ \\
Penyelesaian tugas kelistrikan dan rangkaian DC & 6 & Gaya Coulomb, potensial, kapasitansi, serta arus-tegangan rangkaian seri/paralel dihitung benar dengan hukum Ohm/Kirchhoff dan satuan konsisten. & Sebagian besar perhitungan benar, tetapi ada kesalahan pada analisis rangkaian atau satuan. & Perhitungan kelistrikan salah konsep atau analisis rangkaian tidak dapat diikuti. \\}

\assessmentblock{Kuis Pemahaman Konsep Fisika (Kuis 1 dan Kuis 2)}{Kuis (pilihan ganda dan/atau esai, closed book)}{SCPMK0902-00901, SCPMK0902-00903}{Kuis 1 (pekan 4) menguji materi pekan 1--3: hakikat fisika, metode ilmiah, dan pengukuran. Kuis 2 (pekan 12) menguji materi pekan 10--11: usaha, energi, impuls, dan momentum.}{Menjawab soal pilihan ganda dan/atau esai secara mandiri dalam waktu terbatas tanpa membuka catatan.}{Lembar jawaban kuis.}{Ketepatan jawaban Kuis 1: konsep dasar, metode ilmiah, dan pengukuran & 10 & Jawaban benar minimal 75\% sesuai kunci; langkah dan alasan pada soal esai tepat dan menggunakan konsep yang benar. & Jawaban benar 50--74\%; sebagian langkah esai kurang lengkap atau kurang tepat. & Jawaban benar di bawah 50\% atau langkah penyelesaian tidak sesuai konsep. \\ \\
Ketepatan jawaban Kuis 2: usaha, energi, impuls, dan momentum & 10 & Jawaban benar minimal 75\% sesuai kunci; penerapan hukum kekekalan energi/momentum pada soal esai tepat. & Jawaban benar 50--74\%; penerapan hukum kekekalan sebagian keliru. & Jawaban benar di bawah 50\% atau hukum kekekalan diterapkan salah. \\}

\assessmentblock{UTS Fisika}{UTS (ujian teori tertulis, closed book)}{SCPMK0902-00901, SCPMK0902-00902}{Mahasiswa mengerjakan ujian tengah semester dengan cakupan materi pekan 1--8: hakikat fisika, metode ilmiah, pengukuran, kinematika, dan dinamika partikel.}{Menjawab soal pilihan ganda dan/atau esai secara mandiri dalam waktu terbatas tanpa membuka catatan.}{Lembar jawaban ujian tengah semester.}{Penguasaan konsep dasar, metode ilmiah, dan pengukuran & 10 & Menjelaskan konsep dasar dan tahapan metode ilmiah dengan benar serta menghitung konversi satuan dan ketidakpastian secara tepat pada minimal 75\% butir soal. & Konsep dasar sebagian besar benar (50--74\% butir), tetapi ada kesalahan konversi atau perhitungan ketidakpastian. & Jawaban benar di bawah 50\% butir atau konsep dasar dan pengukuran banyak keliru. \\ \\
Analisis kinematika gerak lurus dan dua dimensi & 10 & Memilih persamaan gerak yang tepat, menghitung parameter GLB/GLBB dan gerak parabola dengan benar, serta membaca grafik gerak dengan tepat. & Persamaan gerak umumnya tepat, tetapi ada kesalahan perhitungan atau pembacaan grafik. & Persamaan gerak dipilih keliru atau hasil perhitungan sebagian besar salah. \\ \\
Penerapan hukum Newton pada dinamika partikel dan sistem mekanik & 10 & Mengidentifikasi gaya-gaya yang bekerja, menyusun resultan gaya, dan menghitung percepatan/gaya dengan benar termasuk pada kasus sistem mekanik. & Identifikasi gaya umumnya benar, tetapi ada kesalahan resultan atau perhitungan percepatan. & Gaya-gaya tidak teridentifikasi benar atau hukum Newton diterapkan salah. \\}

\assessmentblock{UAS Fisika}{UAS (ujian teori tertulis, closed book)}{SCPMK0902-00903, SCPMK0902-00904}{Mahasiswa mengerjakan ujian akhir semester dengan cakupan materi pekan 1--16 dan penekanan pada usaha-energi, momentum-tumbukan, serta kelistrikan statis dan rangkaian DC.}{Menjawab soal pilihan ganda dan/atau esai secara mandiri dalam waktu terbatas tanpa membuka catatan.}{Lembar jawaban ujian akhir semester.}{Analisis usaha, energi, dan momentum-tumbukan & 8 & Menghitung usaha dan konversi energi dengan benar serta menerapkan hukum kekekalan momentum pada tumbukan elastis/inelastis dengan hasil tepat. & Konsep usaha-energi dan momentum umumnya benar, tetapi ada kesalahan perhitungan atau identifikasi jenis tumbukan. & Hukum kekekalan energi/momentum diterapkan salah atau hasil tidak relevan. \\ \\
Analisis gaya Coulomb, medan, dan potensial listrik & 11 & Menghitung gaya Coulomb, kuat medan, energi potensial, dan beda potensial dengan benar serta menggambarkan garis medan dengan tepat. & Sebagian besar perhitungan benar, tetapi ada kesalahan pada hubungan medan-potensial atau penggambaran medan. & Perhitungan gaya/medan/potensial banyak salah konsep. \\ \\
Analisis kapasitansi dan rangkaian listrik DC & 11 & Menghitung kapasitansi, energi tersimpan, serta arus, tegangan, daya, dan efisiensi rangkaian seri/paralel dengan hukum Ohm/Kirchhoff secara benar. & Perhitungan kapasitor atau rangkaian umumnya benar, tetapi ada kesalahan analisis seri/paralel. & Analisis kapasitor dan rangkaian DC salah konsep atau tidak dapat diikuti. \\}

\begin{tblr}{width=\textwidth,colspec={|Q[l,m,3.2cm]|X[l,m]|},hlines,vlines,hspan=minimal,row{1-3}={bg=headgray},rowsep=1.5pt,colsep=3pt}
Jadwal Pelaksanaan: & Minggu 1--17 sesuai RPS. Setiap evaluasi memiliki RTM dan rubrik multi-kriteria. \\
Lain-Lain Yang Diperlukan: & LMS, lembar kerja/tugas, perangkat lunak simulasi fisika, dan perangkat presentasi. \\
Pustaka: & Mengacu pustaka utama dan pendukung pada bagian RPS. \\
\end{tblr}
